\chapter[Limitations and failure modes]{Limitations and failure modes}
\label{ch:failures}
\chapterpromise{Conditions under which the factorized partition model, segment distribution, numerical approximation, hierarchical assumptions, prediction rule, or empirical comparison is not valid.}

\label{sec:limitations}

The exact results in this volume concern offline Bayesian segmentation on a finite ordered
candidate set. They require finite segment marginal likelihoods and a prior whose contribution can
be represented by segment or transition factors. Nonconjugate calculations require additional
numerical assumptions. This chapter states the corresponding limitations.

\section{Computational and model support}
\paragraph{Sequence length and segment count.}
The dynamic program has worst-case time complexity $\Theta(k_{\max}n^2)$ and working memory at
least $\Theta(k_{\max}n)$ for the principal posterior quantities
(Proposition~\ref{prop:bb-complexity}). It is an offline method. Large sequences may require a
restricted changepoint set, pruning, or an approximate decomposition. Such modifications change the
set of partitions or the summation unless their exactness is established separately; they do not
automatically inherit Theorem~\ref{thm:dp-correctness}.

\paragraph{Candidate changepoints.}
A changepoint outside the candidate set cannot receive posterior probability. This is a support
restriction of the model, not a numerical error. Scientific locations that must remain possible
should be included before fitting.

\paragraph{Nonlocal dependence among segments.}
The recursions assume that, after any necessary state variable is included, the probability of a
partition can be written as products of local contributions. A prior depending on a global pattern
of segment lengths, unrestricted interactions among nonadjacent segments, or labels with
non-Markov dependence may violate this assumption. Some dependence structures can be handled by an
enlarged state space, but the recurrence and complexity must then be derived for that model.

\section{Observation-model misspecification and prior sensitivity}
\paragraph{Misspecified segment distribution.}
Posterior probabilities are conditional on the chosen within-segment distribution, its descriptors,
and its hyperparameters. Heavy-tailed observations under a Gaussian segment model, zero inflation
under a Poisson model, overdispersion under a Binomial model, or unmodeled serial dependence can
distort segment marginal likelihoods and therefore the posterior over partitions. Posterior
predictive diagnostics can reveal discrepancies, but the method does not automatically select the
correct observation family.

\paragraph{Prior sensitivity.}
The posterior for $k$ and the changepoint probabilities can depend materially on the segment-count
prior, segment cohesion $c_u(i,j)$, and boundary hazard $h_u(j)$. This dependence can be stronger on
an irregular grid because different priors assign different probabilities to long physical gaps.
An analysis should report sensitivity to scientifically plausible prior choices rather than treat a
single prior as neutral.

\section{Numerical approximation}
\paragraph{Nonconjugate segment marginal likelihoods.}
Proposition~\ref{prop:stability} assumes a verified bound on the largest log marginal-likelihood
error over every admissible segment. The bound increases with the number of segments, so a
fixed segmentwise error can have a larger effect for large $k$. If the numerical method does not
establish this condition, the dynamic program returns the exact posterior for the approximate score
array but no error bound relative to the intended posterior.

In the archived comparison, the maximum absolute discrepancy for expectation propagation was
approximately $14.5$ relative to the quadrature reference on the examined segment set. This is not
a small-error setting. The associated boundary score is descriptive and does not verify the
stability theorem. A scientifically interpretable nonconjugate analysis should compare numerical
methods on the same admissible segments, examine convergence and tail behavior, and use a more
accurate calculation when the discrepancy is too large for the desired posterior comparison.

\section{Hierarchical assumptions}
\paragraph{Common changepoints across sequences.}
The shared-changepoint model assumes aligned candidate coordinates and conditional independence
across sequences given a common partition and sequence-specific segment parameters. Unmodeled
cross-sequence dependence or genuine sequence-specific changepoints can make the common partition
misleading. The finite-candidate consistency theorem also requires a positive average expected
log-score gap between the true partition and every alternative; one informative sequence does not
satisfy that condition as the number of sequences increases.

\paragraph{Latent-group non-uniqueness.}
The finite latent-group criterion is invariant under permutation of the group labels and may have
duplicate or collapsed templates. Multiple restarts do not prove that a global maximum has been
found. Because the sequence-template scores are not assumed to be normalized component densities,
finite-mixture identifiability theorems do not apply. The archived implementation must also be
corrected so that restart ranking uses the final value in each objective sequence.

\paragraph{Posterior summaries versus joint optimization.}
Marginal changepoint probabilities and the joint MAP partition are different summaries. Selecting
the marginal mode for each ordered changepoint coordinate need not yield a valid or jointly most
probable partition. Applications should report the posterior quantities relevant to the decision
rather than interpreting all returned summaries as interchangeable.

\section{Prediction and evaluation}
\paragraph{Coordinate support and extrapolation.}
Posterior prediction must use the fitted observation family. A coordinate outside the fitted range
requires a pre-specified extrapolation model or rejection. Silent endpoint assignment changes the
predictive model and can yield a finite but invalid score, as demonstrated by the archived
methylation calculation excluded from predictive conclusions.

\paragraph{Changepoint metrics.}
Tolerance-based boundary $F_1$ depends on the tolerance and matching rule. When \BayesBreak MAP
boundaries are used as the reference, a comparator score measures agreement with \BayesBreak, not
accuracy against independent truth. A location error is undefined when no pair is matched and is
reported as \NA. The four real-data analyses do not contain independent changepoint annotations,
so their fitted red lines cannot be interpreted as verified geological, genomic, or financial
events.

\paragraph{Incompatible comparator coordinates.}
The array-CGH comparator flattened a $43\times2215$ fitted matrix to 95,245 values while the
reference changepoints remained on the 2,215-probe axis. Those objects do not define a common
matching problem. The executed values remain recorded, but they are excluded from comparator
conclusions.

\section{Choice of analysis}
\begin{enumerate}[label=(\arabic*),leftmargin=2.3em]
\item For a single sequence with a conjugate segment model, use the exact segment integration and
sum-product or max-sum recursions of Chapters~\ref{ch:block}--\ref{ch:map}.
\item For a nonconjugate segment model, use a stated numerical method and assess its error on the
admissible segment set. Use Proposition~\ref{prop:stability} only when
Assumption~\ref{ass:uniform-block-error} has been established.
\item For aligned sequences expected to share changepoints, use the common-partition model of
Chapter~\ref{ch:shared}; examine sensitivity to sequence dependence and violations of common
structure.
\item For observed groups, fit the shared-changepoint model within each group. For unobserved
groups, use the finite latent-group criterion of Chapter~\ref{ch:templates} and interpret its
allocations and templates according to that criterion, without importing normalized-mixture
claims.
\item For prediction, use the posterior predictive distribution of the fitted observation family
and state whether conditioning is on a selected partition or averaged over partition uncertainty.
\end{enumerate}

\section{Empirical coverage}
The synthetic studies and four real-data fits demonstrate the calculations under their archived
settings. They do not provide a complete fairness-controlled comparison with all relevant
frequentist and Bayesian methods, broad repeated uncertainty intervals, comprehensive
misspecification studies, or large-$n$ scaling. Those studies remain necessary before making broad
comparative or robustness claims.

\begin{figure}[H]
\centering
\resizebox{0.96\textwidth}{!}{\begin{tikzpicture}[node distance=9mm and 12mm]
\node[primarybox,text width=32mm] (factor) {Contiguous partition with locally factorized segment contributions};
\node[primarybox,text width=32mm,right=of factor] (segment) {Finite segment marginal likelihood or stated numerical approximation};
\node[successbox,text width=32mm,right=of segment] (exact) {Exact partition recursion or conditional approximation bound};
\node[neutralbox,text width=34mm,below=14mm of factor] (nonlocal) {Nonlocal dependence among segments or unrestricted labels};
\node[neutralbox,text width=34mm,below=14mm of segment] (uncert) {Uncontrolled segment scores or observation-support mismatch};
\node[secondarybox,text width=34mm,below=14mm of exact] (outside) {Requires a different model, state space, or numerical analysis};
\draw[arrPrimary] (factor) -- (segment);
\draw[arrPrimary] (segment) -- (exact);
\draw[arr,dashed] (nonlocal) -- node[below,font=\scriptsize]{local factorization fails} (outside);
\draw[arr,dashed] (uncert) -- node[below,font=\scriptsize]{posterior interpretation not established} (outside);
\draw[arr,dashed] (factor) -- (nonlocal);
\draw[arr,dashed] (segment) -- (uncert);
\end{tikzpicture}
}
\caption[Scope and failure modes]{Scope of the exact partition calculation. Unsupported candidate
sets, nonlocal dependence, misspecified segment distributions, uncontrolled numerical scores,
undefined extrapolation, and incompatible comparator axes require a revised model or analysis.}
\label{fig:scope-failure-map}
\end{figure}
